\documentclass[letterpaper,12pt]{article}
 \usepackage[utf8]{inputenc}
\usepackage[T1]{fontenc}
% \usepackage[latin9]{inputenc}
\usepackage{amstext}
\usepackage{amsmath}
\usepackage{amssymb}
\usepackage{graphicx}
\usepackage{geometry}
\usepackage{tikz}
\usepackage{esint}

\geometry{
 letterpaper,
%  total={210mm,297mm},
 left=1.25in,
 right=1.25in,
 top=1.5in,
 bottom=0.5in,
 }
 
 \begin{document}
  \begin{picture}(0,0)
\put(114,-10){\includegraphics[width=6cm]{Images/un_logo.eps}}
\end{picture}
\begin{center}
\textbf{
 FACULTAD DE INGENIER\'IA\\
 Vicedecanatura Acad\'emica\\
 POSGRADOS \break
  \newline
 PROPOSAL SUBMISSION
}

\end{center}
\vspace{30pt}

% •Título preliminar
% •Justificación
% •Estado del arte
% •Definición del problema
% •Metodología
% •Cronograma aproximado
% •Sectores de impacto de la investigación
% •Referencias bibliográficas 

%checkboxes
\begin{tabular}{l c l c}
 DOCTORAL THESIS: & \framebox[0.5cm][c]{x} & MASTER THESIS: & 
\framebox[0.5cm][c]{}\\
 MASTER FINAL WORK: & \framebox[0.5cm][c]{} & SPECIALIZATION FINAL WORK: & 
\framebox[0.5cm][c]{}\\ 
\end{tabular}
\vspace{20pt}
 \begin{enumerate}
  \item \textbf{BIDDER:} Robinson Andr\'es Jaque Pirab\'an \qquad\qquad 
\textbf{ID:} 80190790
  \item \textbf{PROGRAM:} Phylosophy Doctoral in Computer Science and Systems 
Engineering
  \item \textbf{ADVISOR:} Fabio Augusto Gonz\'alez Osorio\\
  \textbf{DEPARTMENT:} Computer Science and Industrial Engineering
  \item \textbf{TITLE: Kernel Tensor Factorization }%check
  \item \textbf{AREA:} Computer Science
  \item \textbf{LINE OF RESEARCH:} Machine Learning
  \item COMMENTARY WITH ADVISOR APROVAL
  \vspace{120pt}
% % 14. COMENTARIO CON VISTO BUENO DEL DIRECTOR: (calificar los siguientes 
% aspectos: organización, pertinencia, relevancia y originalidad).
  \item BIDDER SIGNATURE 
%  
%  Robinson Andr\'es Jaque Pirab\'an
% % 15. FIRMA DEL PROPONENTE
  \vspace{60pt}
  \item SIGNATURE OF ADVISOR
  \vspace{60pt}
 % 16. FIRMA DEL DIRECTOR (ASESORES)
 % 17. FECHA  
  \end{enumerate}  
  
%   \item \textbf{BACKGROUND AND JUSTIFICATION:}\\
  
% 5. ANTECEDENTES Y JUSTIFICACIÓN: (Indicar los desarrollos previos, 
% circunstancias y condiciones que llevaron a la conclusión de la necesidad y 
% conveniencia del proyecto)

%  \item \textbf{PROBLEM STATEMENT:}\\
% 6. IDENTIFICACIÓN DEL PROBLEMA: 


\clearpage

%Portada

\clearpage

\tableofcontents

\section{Introduction}

Tensors are multidimensional arraies. i.e. An $N$-way or $N$-order tensor is an 
element of tensor product of $N$ vector spaces, each of which has its own 
coordinate system. A first order tensor is a vector, a second order tensor is a 
matrix, tensors of higher order are called high-order tensors.

The order (ways or modes) of a tensor is the number of dimensions. In figure 

Third-order tensor 


\subsection{Tensor decomposition}


Tensor decomposition originated with Hitchcock in 1927 (Reference of Kolda), and 
the the multi-way model is attribuited to Cattel in 1944 (Reference of Kolda).

Tensor works had attention in 60s with Tucker, Carroll and Chang (Reference of 
 Kolda) and Harshman in 1970 with applications in psychometrics. In 1981 
Appellof and Davidson used tensor decomposition in chemometrics which have been 
an popular field of application of tensor decomposition since then.

In last twenty years tensor decomposition applications have expanded to many 
fields such as signal processing, numberical linear algebra, computer vision, 
numerical analysis, neuroscience, data mining, graph analysis.


% include figure about application fields.

% Extend applications descriptions and cite paper and survey Specific for each 
% field.


We suggest to readers to refer to (Kolda, 2009) for an exhaustive a detailed 
review of fundamental decomposition methods.



\subsubsection{Formulation for tensor completion}

Following Ji Lu et. al notation

low rank matrix completion

\begin{equation} 
 \begin{split}
  \min_X & \text{ rank}(X)\\
  \text{s.t. } & X_\Omega = M_\Omega
 \end{split}
\end{equation}
where $\Omega$ is an index set, then $X_\Omega$ is coping entries of $X$ in the 
indexes $\Omega$ and missed entries $\hat{\Omega}$ would be $0$

The missing entries in $X$ are determined in order to minimize the matrix $X$ 
rank. i.e. a non convex optimization problem since rank is nonconvex.

Frequently, trace norm (or nuclear norm) $||\cdot ||_*$ is used to approximate 
the rank of matrices.
Trace norm is the tighest convex envelop for the matrices rank.

\begin{equation} 
 \begin{split}
  \min_X & ||X||_*\\
  \text{s.t. } & X_\Omega = M_\Omega
 \end{split}
\end{equation}

Since tensor is a generalization of the matrix concept, we generalize the 
optimization problem as

\begin{equation} 
 \begin{split}
  \min_{\mathcal{X}} & ||\mathcal{X}||_*\\
  \text{s.t. } & \mathcal{X}_\Omega = \mathcal{T}_\Omega
 \end{split}
\end{equation}

Where $\mathcal{X}$ and $\mathcal{T}$ are $n$-order tensors with identical size.



\subsection{Tensor probability}

Given a sample set


\subsection{Kernel matrix}




%  \item \textbf{GENERAL AND SPECIFIC OBJECTIVES} 
% 7. OBJETIVO GENERAL Y OBJETIVOS ESPECÍFICOS: (Evaluables)

% \subsection*{General objective}
% 
% 
% 
% \subsection*{Specific objectives}
% 
% \item \textbf{METHODOLOGY}\\ 
% % 8. METODOLOGÍA: (Descripción del procedimiento a emplear para obtener los 
% objetivos propuestos).
% % 
% 
% \end{enumerate}
% 
% %  \item \textbf{ACTIVITIES} \\
% % 9. ACTIVIDADES A DESARROLLAR:
% 
% 
% Project activities and results are structured within \textit{Work Packages} 
% (WP) in terms of \textit{Tasks} and \textit{Deliverables}, as described below.
% 
% \subsection*{WP0: PROPOSSAL}
% 
% This WP will assure a suitable scope of propossal due to evaluation of risk 
% and feasibility. In addition, activities and schedule will be delimited in 
% concordance with objectives stablished.
% 
% \subsubsection*{TASKS}
% \begin{itemize}
%  \item T0.1: Propossal review
%  \item T0.2: Literature review.
% \end{itemize}
% 
% 
% \subsubsection*{DELIVERABLES}
% \begin{itemize}
%  \item D0.1: A survey paper reporting the literature review.
%  \item D0.2: Final propossal
% \end{itemize}
% 
% 
% \subsection*{WP1: MANAGEMENT }
% 
% This WP will ensure effective planning, implementation, coordination and 
achievement of the project activities, including successful completion of the 
tasks and timely production of deliverables.
% 
% \subsubsection*{TASKS }
% \begin{itemize}
% \item T1.1: ADMINISTRATION AND MANAGEMENT. Will take care of financial and 
resource accounting. It will also be in charge of managing the relations with 
collaborating institutions and administrative bodies within.
% \item T1.2: PROJECT FOLLOW UP AND QUALITY CONTROL. Will act as the final 
stage before delivery hand over to ensure compliance and coherence. Also, it 
will follow up project progress anticipating corrective actions and assessing 
risk mitigation actions.
% \end{itemize}
% 
% \subsubsection*{DELIVERABLES}
% \begin{itemize}
% \item D1.1: Periodic technical project reports.
% % \item D1.2: Periodic financial project reports
% \end{itemize}
% 
% \subsection*{WP2: MEASURING ROBUSTNESS}
% 
% 
% \subsubsection*{TASKS }
% 
% \subsubsection*{DELIVERABLES }
% 
% \subsection*{WP3: DESIGN OF ROBUST MATRIX FACTORIZATION ALGORITHMS}
% 
% \subsubsection*{TASKS }
% 
% \subsubsection*{DELIVERABLES }
% 
% \subsection*{WP4: PERFORMANCE EVALUATION}
% 
% 
% \subsubsection*{TASKS }
% 
% \subsubsection*{DELIVERABLES }
% 
% \subsection*{WP5: DISSERTATION}
% \subsubsection*{TASKS}
% \subsubsection*{DELIVERABLES }
% 
% \clearpage
% 
% %  \item \textbf{SCHEDULE}\\
% % 10. CRONOGRAMA:
% 
% 
% %  \item \textbf{BIBLIOGRAPHY:}\\
% % 11. BIBLIOGRAFÍA BÁSICA:
% 
% \begin{thebibliography}{References}
% 
% 
% 
% \end{thebibliography}
% 
% 
% %  \item \textbf{EQUIPMENT AND SUPPLIES}\\
% % 12. RECURSOS FÍSICOS: (Especificar la disponibilidad y adjuntar carta de 
% compromiso de la dependencia o empresa cuando sea necesario).
% A work station with a computer with internet and bibliographic material.
% 
% 
% %  \item \textbf{COSTS AND FINANCIAL SOURCES}\\
% % 13. COSTOS DEL TRABAJO Y FUENTES DE FINANCIACIÓN:
% 
% \begin{table}[!h]
% \centering
% \begin{tabular}{l l r}
% \hline
%  \textbf{Concept} & \textbf{Source} & \textbf{Total cost}* \\
%  \hline
%  Salary researcher & Researcher & \$144 \\
%  Advisor salary & Universidad Nacional de Colombia & \$80\\
%  Conferences and events & Universidad Nacional de Colombia/Researcher & \$10 
\\
%  \hline
%  Notebook & Researcher & \$2\\
%  Desktop computer & Research group MindLAB & \$1.6\\
%  High performance computer & Research group MindLAB & \$20\\
%  \hline
%  Office supplies & Researcher & \$1.5\\
%  Bibliography & Universidad Nacional de Colombia & \$1\\
%  \hline
%  \textbf{Total} & & \textbf{\$260.1}
% \end{tabular} 
% \end{table}
% \textit{*Millions of colombian pesos (COP)}
% 
% 
% 
% 
% %  \item COMMENTARY WITH ADVISOR APROVAL
% % 14. COMENTARIO CON VISTO BUENO DEL DIRECTOR: (calificar los siguientes 
% aspectos: organización, pertinencia, relevancia y originalidad).
% %  \item BIDDER SIGNATURE 
%  
%  Robinson Andr\'es Jaque Pirab\'an
% % 15. FIRMA DEL PROPONENTE
% %  \item SIGNATURE OF ADVISOR
% % 16. FIRMA DEL DIRECTOR (ASESORES)
% % 17. FECHA
%  
% % \end{enumerate}
% 
% 
% 
% 


\end{document}
